% Wymaga: booktabs. Nie wymaga pakietu float.
\makeatletter
\@ifundefined{tablehere}{\newenvironment{tablehere}{\par\vskip\intextsep\def\@captype{table}\begin{center}}{\end{center}\par\vskip\intextsep}}{}
\makeatother

\subsection{Optymalizacja mediów}
\subsubsection{Opis optymalizacji}

W pierwotnej wersji serwisu zdjęcia do wpisów były ładowane w oryginalnej postaci
(wielomegapikselowe pliki JPEG), a tło strony stanowił film \texttt{1.webm}
o rozmiarze ok.\ 549\,MiB - był to plik 4k bezporednio pobrany z platformy YouTube i nieobrobiony. Tak duże pliki dominowały nad resztą zasobów
(HTML, CSS, JavaScript), wydłużały czas ładowania strony i zużywały transfer
użytkowników oraz serwera. Podczas optymalizacji film przeskalowano, rozmyto i silnie
skompresowano kodekiem VP9, a zdjęcia zmniejszono i zamieniono z JPEG na WebP.
W pliku \texttt{server/server.js} zmieniono ścieżki zdjęć z \texttt{/img/N.jpg}
na \texttt{/img/N.webp}.

\subsubsection{Konfiguracja i parametry}

Film przetworzono używajc narzędzia FFMPEG, poleceniem:

\begin{verbatim}
ffmpeg -y -i 1.webm \
  -vf "scale=1080:-2:flags=lanczos,fps=24,gblur=sigma=5" \
  -an -c:v libvpx-vp9 -b:v 0 -crf 46 -row-mt 1 2.webm
\end{verbatim}

\begin{tablehere}
\centering
\begin{tabular}{l p{9cm}}
\toprule
\textbf{Parametr} & \textbf{Znaczenie} \\
\midrule
\texttt{scale=1080:-2} & szerokość 1080\,px, wysokość wyliczana z zachowaniem proporcji
  (liczba parzysta, której wymaga koder) \\
\texttt{flags=lanczos} & algorytm Lanczosa przy skalowaniu -- ostrzejszy obraz niż
  domyślny bikubiczny \\
\texttt{fps=24} & ograniczenie liczby klatek do 24 na sekundę \\
\texttt{gblur=sigma=5} & rozmycie Gaussa ($\sigma=5$); film jest tylko
  rozmytym tłem, więc utrata detali jest niewidoczna, a koder potrzebuje
  znacznie mniej danych \\
\texttt{-an} & usunięcie ścieżki dźwiękowej (tło jest niemym wideo) \\
\texttt{-c:v libvpx-vp9} & kodek wideo VP9 w kontenerze WebM \\
\texttt{-b:v 0} & brak docelowej przepływności -- w połączeniu z \texttt{-crf}
  włącza tryb stałej jakości \\
\texttt{-crf 46} & wysoki współczynnik kompresji (skala VP9: 0--63, im więcej,
  tym mniejszy plik i niższa jakość) \\
\texttt{-row-mt 1} & wielowątkowe kodowanie wierszy bloków -- przyspiesza
  kodowanie, nie wpływa na rozmiar \\
\bottomrule
\end{tabular}
\caption{Parametry konwersji filmu}
\end{tablehere}

Zdjęcia konwertowano z użyciem narzędzia FFMPEG, poleceniem:

\begin{verbatim}
ffmpeg -i img/1.jpg -vf "scale=-2:720" \
  -c:v libwebp -quality 60 img/1.webp
\end{verbatim}

\begin{tablehere}
\centering
\begin{tabular}{l p{9cm}}
\toprule
\textbf{Parametr} & \textbf{Znaczenie} \\
\midrule
\texttt{scale=-2:720} & wysokość 720\,px, szerokość wyliczana automatycznie
  z zachowaniem proporcji (podzielna przez 2) \\
\texttt{-c:v libwebp} & kodowanie do formatu WebP \\
\texttt{-quality 60} & jakość kompresji stratnej w skali 0--100 \\
\bottomrule
\end{tabular}
\caption{Parametry konwersji zdjęć}
\end{tablehere}

Po konwersji zdjęcia mają rozdzielczość od $476\times720$ do $1188\times720$\,px,
a film $1080\times608$\,px w kodeku VP9 przy 24\,kl./s.

\subsubsection{Wyniki pomiarów}

(dla filmu -- wskaźnika Git LFS w commicie \texttt{3c11684}).

\begin{tablehere}
\centering
\begin{tabular}{l r r r r}
\toprule
\textbf{Plik} & \textbf{Przed [kB]} & \textbf{Po [kB]} & \textbf{Redukcja} & \textbf{Krotność} \\
\midrule
Zdjęcia (5 plików) & 13\,704 & 528 & 96,1\% & 26,0$\times$ \\
\texttt{video/1.webm} & 575\,451 & 2\,312 & 99,6\% & 248,8$\times$ \\
\midrule
\textbf{Razem} & \textbf{589\,155} & \textbf{2\,840} & \textbf{99,5\%} & \textbf{207,4$\times$} \\
\bottomrule
\end{tabular}
\caption{Rozmiary mediów przed i po optymalizacji (1\,kB = 1000\,B)}
\end{tablehere}

\begin{tablehere}
\centering
\begin{tabular}{l c c c c c}
\toprule
\textbf{Metryka} & \textbf{PSI mobile} & \textbf{PSI desktop} & \textbf{GTmetrix} & \textbf{Zmiana vs. stan wyjściowy} & \textbf{Uwagi} \\
\midrule
LCP [s]        & [—] & [—] & [—] & [—] & [—] \\
TBT [ms]       & [—] & [—] & [—] & [—] & [—] \\
CLS            & [—] & [—] & [—] & [—] & [—] \\
FCP [s]        & [—] & [—] & [—] & [—] & [—] \\
Speed Index [s]& [—] & [—] & [—] & [—] & [—] \\
\midrule
Wynik wydajności (0–100) & [—] & [—] & [—] & [—] & [—] \\
Rozmiar strony [MB]      & —   & —   & [—] & [—] & [—] \\
Liczba żądań HTTP        & —   & —   & [—] & [—] & [—] \\
\bottomrule
\end{tabular}
\caption{Wyniki pomiarów dla optymalizacji mediów}
\end{tablehere}

\subsubsection{Wnioski}

Łączny rozmiar mediów spadł z około 561,9\,MiB do 2,7\,MiB. Zdjęcia zmniejszyły się
średnio 26-krotnie (z 13,1\,MiB do 0,5\,MiB), a film niemal 250-krotnie
(z 548,8\,MiB do 2,2\,MiB).

Największy zysk dało przetworzenie filmu: ponieważ pełni on rolę rozmytego tła,
można było bez zauważalnej straty jakości zmniejszyć rozdzielczość, liczbę klatek,
usunąć dźwięk i zastosować mocną kompresję. W przypadku zdjęć wystarczyło
ograniczenie wysokości do 720\,px oraz zamiana JPEG na WebP o jakości 60.
Kosztem jest nieodwracalna utrata szczegółów, dlatego oryginały warto
przechowywać poza ścieżką serwowaną użytkownikom.
